\section{Published-loss adaptation, perception and reference cost}
\label{sec:video-evidence}
\paragraph{A controlled published-loss reference.}
We adapt the official 4,316,259-parameter StableLLVE U-Net
and deformation-consistency loss form~\cite{stablellve} to the fixed mixed
task. Both arms use the same public clean initialization, source windows,
BGR ordering and 8,000 updates per seed. The spatial arm minimizes L1;
the second adds $20\lVert F(Wx)-WF(x)\rVert_1$, retaining both prediction
paths' gradients. Here $W$ uses training-GT motion as a deformation field
applied to the same noisy source image. It does not create a separately
noised adjacent observation. Compared with the paper, this changes CMP
motion synthesis, crop size, initialization and training schedule. The
original warp coordinate convention and CUDA backward nondeterminism
are disclosed. This is a fixed-budget task adaptation, not a full
reproduction, a convergence claim or a new method.

\begin{table}[t]
\centering\scriptsize\setlength{\tabcolsep}{4pt}
\caption{Complete StableLLVE network/loss task adaptation. Every arm starts from the same public clean checkpoint and receives 8,000 updates per seed. The official L1 and weight-20 commutator form are retained, but training flow source, crop, initialization and schedule differ from the paper. Raw has no history/future input; rec. adds classical motion/history. JOD uses the fixed nominal display conditions in Table~\ref{tab:perceptual}. This is neither a new method nor a full-paper reproduction.}
\label{tab:stablellve-adapted}
\begin{tabular}{llrrrrrr}
\toprule
Loss/output & Seed & SYN PSNR & SYN $E_{\rm temp}$ & SYN JOD & ND PSNR & ND $E_{\rm temp}$ & ND JOD\\
\midrule
L1 raw & 0 & 28.087 & 7.619 & 7.657 & 28.732 & 3.235 & 8.432\\
L1 raw & 1 & 28.033 & 7.622 & 7.639 & 28.420 & 3.697 & 8.208\\
L1 + rec. & 0 & 27.590 & 5.493 & 7.632 & 28.582 & 2.437 & 8.430\\
L1 + rec. & 1 & 27.547 & 5.497 & 7.613 & 28.578 & 2.741 & 8.247\\
L1 + commute20 raw & 0 & 25.368 & 7.743 & 6.992 & 25.853 & 2.905 & 7.765\\
L1 + commute20 raw & 1 & 25.246 & 7.875 & 6.972 & 25.828 & 2.905 & 7.676\\
L1 + commute20 + rec. & 0 & 25.204 & 5.789 & 6.912 & 25.632 & 2.353 & 7.651\\
L1 + commute20 + rec. & 1 & 25.147 & 5.859 & 6.895 & 25.579 & 2.344 & 7.562\\
\bottomrule
\end{tabular}
\end{table}

Table~\ref{tab:stablellve-adapted} contains only the four final 8,000-update
models. Relative to its same-network spatial control, the commutator arm
loses 2.59--2.88\,dB raw PSNR and 0.53--0.67 predicted JOD across domains
and seeds. SYN temporal error increases by 1.63--3.32\%; paired-real
error decreases by 10.19--21.44\%, with lower spatial and perceptual
quality. Applying the same recursive reference does not remove this
trade-off. This fixed loss setting is stopped without a local weight,
geometry, learning-rate, crop or training-length search. It does not
invalidate the original paper or all equivariance objectives.

The L1 control is a useful stronger spatial reference: SYN raw PSNR is
28.03--28.09\,dB, and its recursive temporal error is 5.49--5.50 versus
6.39--6.55 for the MobileIE recursive reference. However, its raw temporal
error remains worse than that strong recursive reference. On paired-real
seed 1, even its recursive error is 2.74 versus MobileIE's 2.53. Capacity,
loss, pretraining and compute differ together; their separate effects
cannot be inferred from this comparison. All 20 preassigned static
contacts were inspected, including native center crops. These contacts
are not full-video playback or human flicker judgments.

\paragraph{Matching the spatial objective does not close the gap.}
\begin{table}[t]
\centering\scriptsize\setlength{\tabcolsep}{4pt}
\caption{Matched original MobileIE objective settings, final 8,000 updates only. Both objectives use the same source-window schedule, expanded training parameterization and two-stage Adam budget; L1/MSE gradient scale interacts with clipping. Every row deploys 4,075 parameters. Rec. adds the same classical motion/history reference. SYN comprises three source videos under two noise conditions; ND comprises four reused development clips. JOD uses the fixed assumptions in Table~\ref{tab:perceptual}. No objective search, human study, handset or independent final-test claim.}
\label{tab:matched-objective}
\begin{tabular}{llrrrrrr}
\toprule
Objective/output & Seed & SYN PSNR & SYN $E_{\rm temp}$ & SYN JOD & ND PSNR & ND $E_{\rm temp}$ & ND JOD\\
\midrule
MSE raw & 0 & 26.475 & 10.102 & 7.040 & 27.252 & 3.729 & 8.414\\
MSE raw & 1 & 26.456 & 9.699 & 7.111 & 27.004 & 3.577 & 8.552\\
MSE + rec. & 0 & 26.800 & 6.546 & 7.424 & 27.292 & 2.605 & 8.373\\
MSE + rec. & 1 & 26.613 & 6.393 & 7.409 & 26.923 & 2.534 & 8.512\\
L1 raw & 0 & 26.301 & 9.836 & 7.083 & 26.227 & 3.685 & 7.936\\
L1 raw & 1 & 26.492 & 9.351 & 7.104 & 26.633 & 3.474 & 8.428\\
L1 + rec. & 0 & 26.545 & 6.443 & 7.408 & 26.189 & 2.600 & 7.892\\
L1 + rec. & 1 & 26.531 & 6.279 & 7.354 & 26.566 & 2.512 & 8.382\\
\bottomrule
\end{tabular}
\end{table}

We separately train original MobileIE with spatial L1 under the same
source-window schedule and two-stage 8,000-update budget as its existing
MSE reference (Table~\ref{tab:matched-objective}). The original expanded
training graph and initialization are retained; both deploy the same
4,075-parameter graph. The comparison concerns the full objective setting:
different gradient scales interact with the shared gradient-clipping rule.
Only the preassigned final weights are evaluated.

L1 does not give a consistent joint improvement. Raw temporal error falls
by 1.18--3.59\%, while paired-real raw PSNR falls by 0.37--1.03\,dB and
predicted JOD by 0.12--0.48. Recursive temporal error falls by only
0.22--1.80\%, again with quality costs. SYN seed 1 raw gains just
0.036\,dB while predicted JOD decreases slightly. This fixed reference
is retained without an objective-mixture, learning-rate or seed search.
All 20 prescribed contacts were inspected, without full-video playback
or a human temporal study.

With both spatial references trained using L1, the U-Net still has
1.54--1.79\,dB higher SYN raw PSNR and 1.79--2.51\,dB higher paired-real
raw PSNR. Thus the objective difference alone does not account for the
observed gap, but architecture, capacity, pretraining and optimization
remain coupled. This comparison does not identify capacity as the cause
or establish a new temporal method.

\paragraph{Perceptual predictions qualify the temporal proxy.}
\begin{table}[t]
\centering\small
\caption{Full-reference ColorVideoVDP v0.5.7 predictions for the existing references, with both adapted seeds. The original StableLLVE clean checkpoint has zero project updates and no training seed. Nominal 30\,fps, original-sized SDR arrays, standard FHD display, 37.84 pixels/degree, 200\,cd/m$^2$ peak and replicate temporal padding are assumed. These JOD predictions are not measured handset viewing, capture timing, human MOS or a separately calibrated flicker score.}
\label{tab:perceptual}
\begin{tabular}{llrr}
\toprule
Reference/output & Seed & SYN JOD & ND JOD\\
\midrule
MobileIE raw & 0 & 7.040 & 8.414\\
MobileIE raw & 1 & 7.111 & 8.552\\
MobileIE + rec. & 0 & 7.424 & 8.373\\
MobileIE + rec. & 1 & 7.409 & 8.512\\
Nano raw & 0 & 6.742 & 8.169\\
Nano raw & 1 & 6.872 & 8.141\\
Nano + rec. & 0 & 7.118 & 8.114\\
Nano + rec. & 1 & 7.178 & 8.080\\
FastDVDnet window & 0 & 7.432 & 7.680\\
FastDVDnet window & 1 & 6.304 & 7.257\\
StableLLVE clean raw & -- & 6.262 & 6.802\\
StableLLVE clean + rec. & -- & 6.266 & 6.632\\
\bottomrule
\end{tabular}
\end{table}

We evaluate ColorVideoVDP~\cite{colorvideovdp} on the complete cached
sequences with one fixed display configuration
(Table~\ref{tab:perceptual}). The combined collection has 220
sequence-condition rows and 9,856 comparison frames, reusing 448 input
observations; ten GT identity controls return exactly 10. These are
repeated conditions, not independent video samples. In particular, the
six SYN conditions are three source videos under two noise conditions.
No enhancement network is rerun to obtain these scores.

For MobileIE on SYN, recursion improves predicted JOD by 0.30--0.38
alongside lower temporal error. On paired-real, the mean JOD changes
are about $-0.04$ despite roughly 29--30\% lower temporal error. The
negative change is concentrated in video 15, while the other three
sources have positive changes for both seeds. These small means are
not evidence of statistically significant or human-perceived degradation.
The spatially adapted U-Net also shows that lower temporal error need
not improve JOD: on SYN its recursive JOD is about 0.026 below raw.
Predicted perceptual quality, spatial fidelity and the temporal residual
therefore need joint reporting. A dedicated human temporal study remains
necessary; no display-condition grid is used to select favorable scores.

\paragraph{Efficiency remains a separate empirical requirement.}
\begin{table}[t]
\centering\small
\caption{Reference deployment conditions at $1080\times1920$, batch one, per output frame. Conv GMAC counts only executed Conv2d multiply--accumulates; it omits bias, resizing, priors, pooling, activation, motion, warping and memory work. FastDVDnet uses the original uncached three first-stage and one second-stage calls. Rec. additionally retains two FP32 RGB frames plus motion/workspace; raw spatial networks have no temporal state. Identical StableLLVE graphs serve its public and adapted weights. No row is a measured handset latency or complete pipeline FLOP claim.}
\label{tab:reference-budget}
\begin{tabular}{lrrrrl}
\toprule
Network/output & Params & Conv GMAC & Future & Flow & Handset ms\\
\midrule
MobileIE raw & 4075 & 7.938 & 0 & no & --\\
MobileIE + rec. & 4075 & 7.938 & 0 & yes & --\\
Nano raw & 696 & 1.313 & 0 & no & --\\
Nano + rec. & 696 & 1.313 & 0 & yes & --\\
FastDVDnet window & 2479096 & 1319.208 & 2 & no & --\\
StableLLVE raw & 4316259 & 317.240 & 0 & no & --\\
StableLLVE + rec. & 4316259 & 317.240 & 0 & yes & --\\
\bottomrule
\end{tabular}
\end{table}

The common Conv2d convention in Table~\ref{tab:reference-budget} gives
317.24\,GMAC for StableLLVE versus 7.94 for MobileIE at 1080p. This
approximately forty-fold convolution-budget difference cannot be
translated into a phone-latency ratio. The count excludes motion,
warping, point operations and memory traffic, and the largest individual
convolution output is not peak memory. Neither the 4,075-parameter
starting point nor zero added inference cost is an imposed hard limit.
A defensible deployment claim requires actual full-pipeline phone
measurements for a candidate that first demonstrates joint quality and
stability gains. No such new candidate is established by these references.
